% C.3bq: independently audited restricted theorem; no universal vanishing claim.
\paragraph{C.3bq. Edge and triple-point correction.}
We retain the threefold exceptional data of C.3bm. Thus $k$ is
algebraically closed of characteristic $p>0$, the projective SNC surface
$E=\bigcup_iE_i$ is a divisor in the regular threefold $Z$, and
$D=\sum_i a_iE_i$, with positive integers $a_i$, satisfies
$\mathcal O_E(-D)$ ample. We index every connected component of a double
curve separately, writing $C_e$ for the edge $e$, and every triple point
separately, writing $v$ for a face. There are no fourfold intersections.
The resulting finite dual complex is denoted by $\Delta$. Its rational
cochain groups have the ordered-component orientations and differential
\[
 C^0(\Delta,\mathbf Q)\xrightarrow{\delta^0}
 C^1(\Delta,\mathbf Q)\xrightarrow{\delta^1}
 C^2(\Delta,\mathbf Q).
\]
If $v$ lies on $C_e$, let $q_{e,v}\in C_e(k)$ be that point and let
$\epsilon_{e,v}\in\{1,-1\}$ be its incidence sign. For
$z=(z_v)_v\in H_2(\Delta,\mathbf Q)=\ker(\partial_2)$ define
\begin{equation}
 \mathcal P_e(z)=
 \left[\sum_{v\supset e}\epsilon_{e,v}z_vq_{e,v}\right]
       \in\operatorname{Pic}^0(C_e)(k)\otimes\mathbf Q.
                                                        \tag{ET.1}
\end{equation}
The divisor has degree zero because $\partial_2z=0$ at the edge $e$.
This definition does not require a base point on $C_e$.

Put $S(E)=F_\gamma^2K_0(E)_{\mathbf Q}$ and use the actual proper-curve
relation group $J_K(E)$ of EP. Define
\begin{equation}
 \varphi:S(E)\longrightarrow
          \bigoplus_i\mathrm{CH}_0(E_i)_{\mathbf Q},\qquad
 s\longmapsto\bigl(\operatorname{ch}_2(s|_{E_i})\bigr)_i.
                                                        \tag{ET.2}
\end{equation}
We prove the following two separately scoped assertions.

\emph{First assertion.} If $H^2(\Delta,\mathbf Q)=0$, then
\begin{equation}
                         \ker\varphi\subset J_K(E).     \tag{ET.3}
\end{equation}
\emph{Period criterion.} More generally, (ET.3) holds if
\begin{equation}
 \mathcal P_e(z)=0\quad\text{for every edge $e$ and every }
                 z\in H_2(\Delta,\mathbf Q).             \tag{ET.4}
\end{equation}
The first assertion is a special case of the second, but its proof will
also show directly how the zero-endpoint corrections work.

Together with C.3bo, either hypothesis implies
\begin{equation}
 B^2(Z\setminus E)_{\mathbf Q}=0                         \tag{ET.5}
\end{equation}
if every component has zero rational Albanese kernel, namely if
$\mathrm{CH}_0(E_i)^0_{\mathbf Q}\to
\operatorname{Alb}(E_i)(k)\otimes\mathbf Q$ is injective for every $i$.
In particular this applies when all components are rational surfaces.
No assertion is made that (ET.4), or the component condition, holds for
every exceptional SNC surface.

\begin{proof}
\emph{1. We keep the homotopy descent data, not only the groups.}
For an actual smooth projective parameter scheme $P$, set
\[
 \mathcal C_P^0=\bigoplus_i K(E_i\times P)_{\mathbf Q},\qquad
 \mathcal C_P^1=\bigoplus_e K(C_e\times P)_{\mathbf Q},\qquad
 \mathcal C_P^2=\bigoplus_v K(P)_{\mathbf Q}.
\]
These are nonconnective perfect-$K$ spectra. Face maps are the actual
derived pullbacks. Their ordered alternating sums, together with the
canonical coherences of repeated pullback, give a two-stage homotopy
totalization. Explicitly put
\begin{equation}
 \mathcal T_P=\operatorname{fib}
       (\mathcal C_P^1\longrightarrow\mathcal C_P^2),
 \qquad
 K(E\times P)_{\mathbf Q}\simeq
       \operatorname{fib}(\mathcal C_P^0\longrightarrow\mathcal T_P).
                                                        \tag{ET.6}
\end{equation}
The map from $\mathcal C_P^0$ uses the canonical nullhomotopy of the
double alternating restriction, supplied by the commutative geometric
closed-cover diagram. We do not choose a new nullhomotopy after passing
to its homotopy groups.

Here is the descent input and its exact scope. For a binary closed union
in characteristic $p$, the affine rings form a Milnor square.
\citet[Theorem C, p.~30]{rootGeisserHesselholt2011} makes its bi-relative
groups zero after inverting $p$ in every integer degree. Zariski descent
\citep[Theorem 8.1]{D_ThomasonTrobaugh1990} globalizes the square.
Finite induction on the components and intersections gives this homotopy
limit. Theorem A of the same primary removes any nilpotent reductions
inserted in the induction.
Grouping this finite homotopy limit by double and triple strata gives
(ET.6). Equivalently it is the normalized totalization of the actual
closed-cover diagram; no strata of degree at least three occur.
This construction is natural in $P$, including its point sections.
It is the descent construction used in TS. The argument needs no cellular
splitting and does not assert that individual component families lift.

In particular we have the exact sequences
\begin{align}
 \bigoplus_vK_2(P)_{\mathbf Q}&\longrightarrow\pi_1\mathcal T_P
   \longrightarrow\bigoplus_eK_1(C_e\times P)_{\mathbf Q}
   \longrightarrow\bigoplus_vK_1(P)_{\mathbf Q},        \tag{ET.7}\\
 \pi_1\mathcal T_P&\longrightarrow K_0(E\times P)_{\mathbf Q}
       \longrightarrow\bigoplus_iK_0(E_i\times P)_{\mathbf Q}.
                                                        \tag{ET.8}
\end{align}
The composite of the first arrows in (ET.7)--(ET.8) is the canonical
double-boundary map. Changing a lift in (ET.7) changes the resulting
class only by this double boundary. Thus all coherent-lift ambiguity
is retained explicitly.

\emph{2. Every vertex ambiguity is an actual proper-family relation.}
For $P=\operatorname{Spec}k$, the image of
\[
                  \bigoplus_vK_2(k)_{\mathbf Q}
                      \longrightarrow K_0(E)_{\mathbf Q}
\]
is contained in $J_K(E)$. We recall the actual proper-curve input.
For $a,b\in k^*$ with $a,b\ne1$, set
$c=(a-1)^2(b-1)^2/(ab)$ and normalize a reduced integral component
through $(a,b)$ of
\[
                   (f-1)^2(g-1)^2=cfg
                \quad\text{in }\mathbf P^1\times\mathbf P^1.
\]
There is no vertical or horizontal component. The normalization is a
smooth connected proper curve $P$; its selected component contains the
unique smooth boundary point $(0,1)$, since its $f$-map is surjective.
The symbol $\{f,g\}$ has trivial tame residues at every boundary point.
Localization and Gersten therefore give $\xi\in K_2(P)$ whose two
evaluations are $\{a,b\}$ and zero. At $(0,1)$ one writes
$f=u(g-1)^2$ with $u$ a unit; the Steinberg relation makes the second
evaluation zero. This argument includes characteristic two. The exact
primary inputs are \citet[Section 7, Proposition 3.2, equation (3.3),
and Theorem 5.11]{rootQuillen1973}; the full proper-curve symbol proof
was retained and audited in EL. For its field-generation input we use
Matsumoto's theorem in \citet[Chapter III, Section 6, Theorem 6.1,
chapter p.~48]{frWeibelKbookIII}. The opened source states this theorem
and refers its proof to Milnor, Section 12. Finite sums of symbols
therefore suffice.

Apply the natural double boundary on the whole product to $\xi$ at one
vertex and zero at all other vertices. On the projective threefold
$E\times P$, apply the rational Adams-weight-two projector in $K_0$.
Such a projector exists because $F_\gamma^4K_0(E\times P)=0$, and
the graded Adams eigenvalues $1,2,4,8$ are distinct
\citep[Chapter II, Proposition 4.9 and Remark 8.8.5]
{D_Weibel2013KbookII}. At a point the field $K_2$ class has weight two,
and the double boundary commutes with Adams operations. Here the exact
operation input is the arbitrary-scheme, iterated-homotopy-fibre and
delooping construction in \citet[p.~258, paragraph preceding Lemma 2.1]
{frC3bcLevine1994}. We use its functorial additive Adams operations,
including their agreement with classical operations on $K_0$, without
applying the regular-ambient hypothesis of Lemma 2.1 to the singular
scheme $E$. The projected
family therefore has exactly the prescribed weight-two endpoint class.
This proves the claim. It makes no assertion that a whole class on
$E\times P$ lifts to $Z\times P$.

\emph{3. We remove the degree-unit terms by actual component classes.}
For a smooth proper curve over $k$, localization shows that
$K_1(C_e)_{\mathbf Q}^{(2)}$ is generated by
$u[\mathcal O_c]$, with $u\in k^*$ and $c\in C_e(k)$.
Indeed the generic field $K_1$ has weight one, so a weight-two class is
in the point-pushforward image. Equivalently these are
$u([\mathcal O_{C_e}(c)]-1)$; rational linear combinations allow all
divisor line bundles on $C_e$.

The restriction-degree map
\begin{equation}
 \bigoplus_i\operatorname{Pic}(E_i)_{\mathbf Q}
   \longrightarrow\bigoplus_e\mathbf Q,
 \qquad (L_i)_i\longmapsto
       \bigl(\deg(L_i|_{C_e})-\deg(L_j|_{C_e})\bigr)_{e=(i,j)}
                                                        \tag{ET.9}
\end{equation}
is surjective under the negative ample combination hypothesis. We
recall the numerical proof to avoid turning it into a Picard-restriction
surjectivity claim. A real functional annihilating (ET.9) has
antisymmetric coefficients $t_{ij}$ such that
$\sum_jt_{ij}C_{ij}$ is numerically zero on every component.
Put $u_{ij}=t_{ij}/(a_i a_j)$ and square the normalized relations
$\sum_j a_j u_{ij}C_{ij}=0$, with weight $a_i$ on component $i$.
The normal-bundle interchange on
a transverse double curve, and all transverse triple intersections,
give
\[
 0=-\sum_ea_i a_j\deg((-D)|_{C_e})u_{ij}^2
   -\sum_{v=(i,j,k)}a_i a_j a_k
                          (u_{ij}+u_{jk}+u_{ki})^2.
\]
The first coefficients are strictly positive, so all coefficients
vanish. This proves (ET.9) over $\mathbf Q$. Multiple connected double
curves and multiple triple points are indexed separately throughout.

Now multiply chosen divisor classes in (ET.9) by the constant units
in $K_1(k)$. Their actual classes in $K_1(E_i)$ restrict to the desired
degree-unit edge vectors, possibly together with degree-zero Picard
terms. The image of $\pi_1\mathcal C_k^0$ in
$\pi_1\mathcal T_k$ has zero boundary in (ET.8). We may therefore
subtract these component classes without changing the represented
class in $K_0(E)$, modulo the explicitly retained vertex ambiguity.
We have not asserted that the entire edge vector is a component
restriction: its remaining degree-zero Picard terms are precisely
what the next step treats.

\emph{4. We construct a proper family for the remaining edge vector.}
Take $s\in\ker\varphi$. The rational Chern character on the smooth
surfaces $E_i$ makes its restriction zero in $K_0(E_i)_{\mathbf Q}$.
By (ET.8) it lifts to $\pi_1\mathcal T_k$. Apply the polynomial
$(\psi^2-2)/2$ to this lift. Its boundary remains $s$, because $s$
has weight two. On a smooth proper curve the image of $K_1(C_e)$
in the generic field is $k^*$, split by constant units; its remaining
subgroup is generated by the point pushes above. These two parts have
weights one and two respectively. Thus this polynomial makes the edge
vector belong to $\bigoplus_eK_1(C_e)_{\mathbf Q}^{(2)}$.
We do not assume that $\pi_1\mathcal T_k$ has a semisimple Adams
decomposition. The ambiguity in (ET.7) remains the field $K_2$ term
already treated in Step~2.
After Step~3, linearity and clearing denominators reduce to vectors
of the form
\begin{equation}
               \gamma_e=u([L_e]-1),\qquad
                 L_e\in\operatorname{Pic}^0(C_e)(k),    \tag{ET.10}
\end{equation}
with one fixed unit $u$ and finitely many edges. Finite sums of such
vectors recover the general case.

Choose base points $o_e\in C_e(k)$, put $A=\prod_e\operatorname{Pic}^0(C_e)$,
and choose a smooth connected proper curve $P\subset A$ through $0$
and the tuple $(L_e)_e$. The zero-dimensional case is trivial.
Write $o,t\in P(k)$ for those two points. The normalized Poincare
bundles give actual line bundles $\mathcal L_e$ on $C_e\times P$,
trivial along $o_e\times P$, with endpoint bundles $\mathcal O$ and
$L_e$. Define
\[
       \widetilde\gamma_e=u([\mathcal L_e]-1)
                       \in K_1(C_e\times P)_{\mathbf Q}.
\]
Its edge endpoint vector is zero at $o$ and (ET.10) at $t$.
Its vertex defect is
\begin{equation}
 \eta_v=\sum_{e\subset v}\epsilon_{e,v}\,
             u([M_{e,v}]-1),\qquad
 M_{e,v}=\mathcal L_e|_{q_{e,v}\times P}.                \tag{ET.11}
\end{equation}
Every $\eta_v$ lies in the rational subgroup $N(P)\subset K_1(P)_{\mathbf Q}$
generated by $u'([M]-1)$ for constant units $u'$ and degree-zero line
bundles $M$ on $P$. Every element of $N(P)$ has zero evaluation at
every $k$-point of $P$, since a line bundle on a point has class one.

If $H^2(\Delta,\mathbf Q)=0$, then
$\delta^1\otimes N(P)$ is surjective. Choose
$(\nu_e)_e\in C^1(\Delta,\mathbf Q)\otimes N(P)$ with
$\delta^1\nu=\eta$, and replace
\begin{equation}
       \widetilde\gamma_e\quad\text{by}\quad
       \widetilde\gamma_e-\operatorname{pr}_P^*\nu_e.
                                                        \tag{ET.12}
\end{equation}
The correction restricts to the same $\nu_e$ at every triple point
of that connected edge, and it has zero endpoint fibres. Thus the
corrected vector has zero defect in the last group of (ET.7), while
preserving its two prescribed edge endpoints.

\emph{5. The period criterion gives the same correction.}
Assume (ET.4), without assuming $H^2(\Delta,\mathbf Q)=0$.
For $z\in H_2(\Delta,\mathbf Q)$, pair (ET.11) with $z$ and group
its summands by edge. The normalized Poincare property identifies
the resulting line-bundle combination on $P$, for that edge, with
the pullback of the Poincare evaluation at the degree-zero divisor
in (ET.1). On a smooth curve,
$K_0(P)=\mathbf Z\oplus\operatorname{Pic}(P)$ as groups; consequently
the assignment $M\mapsto[M]-1$ is additive on line bundles, and
torsion line bundles have zero rational class. Hence (ET.4) gives
\[
                       \sum_v z_v\eta_v=0
                          \quad\text{in }N(P).
\]
Finite-dimensional rational linear algebra says that the annihilator
of $\ker\partial_2$ in $C^2(\Delta,\mathbf Q)\otimes N(P)$ is exactly
$\operatorname{im}(\delta^1\otimes N(P))$. Thus the same $\nu_e$ and
the same endpoint-zero correction (ET.12) exist. Notice that this
uses actual Picard classes of the triple-point divisors, not merely
their degrees. Nonzero periods are not silently discarded.

\emph{6. We finish the coherent lift and compare its endpoints.}
Exactness of (ET.7) now lifts the corrected edge vector to
$\pi_1\mathcal T_P$. Its boundary in (ET.8) is an actual class on
$E\times P$. Apply the rational weight-two projector in $K_0$ as
in Step~2 to obtain $\beta\in F_\gamma^2K_0(E\times P)_{\mathbf Q}$.
At the two endpoints its edge vectors are still zero and (ET.10).
The difference between $\beta_t-\beta_o$ and the prescribed class
$s$ is therefore in the double-boundary image of
$\bigoplus_vK_2(k)_{\mathbf Q}$, by (ET.7)--(ET.8) and the component
subtractions of Step~3. Step~2 puts that entire image in $J_K(E)$.
Since $\beta_t-\beta_o$ is itself a defining actual proper-family
relation, $s\in J_K(E)$. This proves (ET.3) in both stated cases.

Finally C.3bo represents every class in the ordinary top quotient
by an element of $S_{00}(E)$. Under the component Albanese-kernel
assumption, all of its component zero-cycles vanish, so it belongs
to $\ker\varphi$. The just-proved inclusion kills it already in
$J_K(E)$ and proves (ET.5).
\end{proof}

\emph{What remains outside the criterion.}
When a triple-point period in (ET.1) is nonzero, the Poincare family
can have a nonzero vertex defect in the cokernel of $\delta^1$.
We have not shown that this defect survives arbitrary endpoint-zero
component families, nor that it gives a nonzero local Chow class.
Conversely we have not found a correction that removes it in general.
The theorem also leaves the component Albanese kernels untouched
unless their vanishing is separately known. These are explicit
limits of the restricted result.
